\clearpage\begingroup\SourceLeaf{694}
\ArithmeticHead{59}{R}{V. ARBEITEN ZUR PRAKTISCHEN MECHANIK.}
\renewcommand{\thefootnote}{\arabic{footnote})}\setcounter{footnote}{121}
\begin{center}\textls[100]{4. Zur Theorie der Wage.}\end{center}
\noindent\hspace*{\parindent}Die Arbeiten an der Regulierung des Hannoverschen Masssystems, die
\textsc{Gauss} zusammen mit \textsc{Schumacher} in den Jahren 1836 bis 1840 durchzuführen
gezwungen war\footnote{Vergl. die \textsc{Gauss}schen amtlichen Berichte darüber in Werke XI, 1, S. 3—15. Ferner den \textsc{Gauss-Schumacher}-Briefwechsel aus den Jahren 1836—1839, Band 2 und 3.}, gaben \textsc{Gauss} Veranlassung, sich mit einer Reihe praktischer
Fragen aus der Theorie der Wägung zu befassen. Die damaligen Mittel erforderten
bei Präzisionsmessungen ganz besondere Vorsichtsmassnahmen, und von den
Schwierigkeiten dieser Messungen und den vielen zu berücksichtigenden Umständen
gibt der Briefwechsel mit \textsc{Schumacher} ein eindrucksvolles Bild. Zwei Fragen
seien hier herausgegriffen, sie betreffen a) ein gegenüber der damals allgemein
üblichen \textsc{Borda}schen Tariermethode abweichendes Wägungsverfahren, b) die
Berichtigung der Schneiden einer Wage.

a) Die auch heute am meisten gebräuchliche Wägungsmethode ist die
auf \textsc{Borda} zurückgehende Tariermethode, bei der eine feste Tara auf der einen
Wagschale bleibt, während auf die andere die zu bestimmende Last bezw.
ein Gewichtssatz gelegt werden, die mit der Tara zu vergleichen sind. \textsc{Gauss}
ersetzte sie durch die Methode der Doppelwägungen, bei der zwei Wägungen
vorgenommen werden, in denen Gewicht und Last zu vertauschen sind\footnote{Brief an \textsc{Schumacher} vom 24. Juli 1836, \textsc{Gauss-Schumacher}-Briefwechsel 3, S. 100. Die \textsc{Gauss}sche Methode war schon vorher von \textsc{Karsten} 1769 entwickelt worden.}.
\textsc{Gauss} erkennt den Vorzug seiner Methode besonders bei oft wiederholten
Wägungen; denn berücksichtigt man nur die zufälligen Beobachtungsfehler,
so ist eine \textsc{Gauss}sche Wägung ebenso genau wie vier \textsc{Borda}sche

b) Bei seinen wiederholten Wägungen entdeckte \textsc{Gauss} als neue Fehlerquelle
der Wage den mangelnden Parallelismus der Schneiden, an denen die
Tragstücke aufsitzen, mit der mittleren Schneide, durch die der Wagebalken
aufsitzt\footnote{\textsc{Gauss-Schumacher}-Briefwechsel, 3, S. 84, 87, 142.}. Sind nämlich die Schneiden nicht parallel, so wird nach dem
Auslösen das Tragstück der Schalen sich nie wieder genau so wie vorher auf
die Schneiden auflegen und daher ein abgeändertes Moment ergeben. Während
bis zu jener Zeit die Schneiden im allgemeinen nur durch Visieren längs der
Schneiden mit blossem Auge nach einem entfernten Objekte berichtigt wurden,
beschäftigt sich nun \textsc{Gauss} eingehender mit der Frage einer genauen Nachprüfung
\par\noindent\hfill 8\textsuperscript{*}
\clearpage\SourceLeaf{695}\ArithmeticHead{60}{L}{H. GEPPERT, ÜBER GAUSS’ ARBEITEN ZUR MECHANIK UND POTENTIALTHEORIE.}
\noindent dieses Schneidenparallelismus. \textsc{Schumacher} bringt in seinen Briefen
vom 7. und 8. Februar 1837\footnote{\textsc{Gauss-Schumacher}-Briefwechsel, 3, S. 147—151.} dafür zwei Methoden in Vorschlag; die erste
besteht darin, dass an jede Schneide ein kleines Fernrohr angesetzt wird, mittels
dessen nach einem entfernten Objekte visiert wird; sie ist wohl praktisch nie
benutzt worden; die zweite wurde von \textsc{Repsold} verwandt und besteht darin,
dass der Wagebalken, der die drei Schneiden trägt, zunächst senkrecht aufgestellt
und dann an die drei nunmehr horizontalen Schneiden eine Libelle
angepresst wird, die deren gegenseitige Neigung zu messen gestattet. Auch
diese Methode hat nicht genügend scharfe Resultate gezeitigt. Deshalb brachte
\textsc{Gauss} zur gleichen Zeit zwei andere Verfahren in Vorschlag; das erste erläutert
er im Briefe vom 12. Februar 1837 an \textsc{Schumacher}\footnote{\textsc{Gauss-Schumacher}-Briefwechsel, 3, S. 155.}. Es ist mit der
heute meist \textsc{Hartig} (1867) zugeschriebenen \textit{Methode mittels Probiergehänge} identisch
und benutzt ein mit zwei Ösen versehenes Tragstück, das verschiedene
Anbringung der Wagschalen gestattet; sind die Schneiden nicht parallel, so
liefern die Wägungen bei Benutzung verschiedener Ösen verschiedene Resultate.
Dieses Verfahren wird als das bequemste in etwas abgeänderter Form
noch heute verwandt. Die zweite \textsc{Gauss}sche Methode wurde in den Göttinger
Gelehrten Anzeigen 1837\footnote{41. Stück vom 13. März 1837, S. 401—405. Werke V, S. 511—513.} veröffentlicht; sie schliesst an einen ebenda erschienenen
Aufsatz von \textsc{W. Weber}\footnote{\textsc{W. Weber}, \textit{De tribus novis librarum construendarum methodis}, ebenda, 22. Stück vom 9. Februar 1837, S. 209—222.} an, in dem sich \textsc{Weber} mit der Aufgabe
beschäftigt hatte, die Schneiden und Pfannen der Wage zu beseitigen
oder die durch sie bedingten Ungenauigkeiten möglichst gering zu gestalten.
Das \textsc{Gauss}sche Korrektionsverfahren besteht darin, dass an eine äussere Schneide
senkrecht zu ihr ein Planspiegel angebracht wird, in dem durch ein Fernrohr
das Bild einer Skala beobachtet wird. Ist die äussere Schneide mit der mittleren
nicht parallel, aber in einer Ebene gelegen, so steigt und sinkt das
Skalenbild im Spiegel, wenn der Wagebalken um die mittlere Schneide
schwingt; sind die Schneiden sogar zu einander windschief, so erleidet das
Bild während des Schwingens auch eine seitliche Verrückung. Gleichzeitig
kann man den Ausschlägen die Grösse der anzubringenden Korrektion entnehmen.
Dieses \textsc{Gauss}sche Verfahren ist auch heute noch in Gebrauch.
\par\medskip\centerline{\rule{.25\textwidth}{.3pt}}
\clearpage\endgroup

\begingroup\NoSourceHead\SourceLeaf{696}
\newcommand{\GCrow}[4]{\par\noindent\hangindent=1.8em\hangafter=1 #1\enspace #2\hfill\mbox{#3}\nobreak\quad #4\par}
\begin{center}\bfseries Inhaltsübersicht.\end{center}
\small
\noindent\hfill Seite.\par
\GCrow{I.}{Nachweis der Erdrotation}{.\ .\ .\ .\ .\ .\ .\ .\ .\ .\ .\ .\ .\ .\ .\ .\ .\ .\ .\ .\ .\ .\ .\ .}{3—16}
\GCrow{\hspace{1.8em}1.}{Fallversuche und ihre Theorie vor \textsc{Gauss}}{.\ .\ .\ .\ .\ .\ .\ .\ .\ .\ .\ .\ .\ .\ .\ .\ .}{3}
\GCrow{\hspace{1.8em}2.}{Die \textsc{Gauss}sche Arbeit}{.\ .\ .\ .\ .\ .\ .\ .\ .\ .\ .\ .\ .\ .\ .\ .\ .\ .\ .\ .\ .\ .\ .\ .}{7}
\GCrow{\hspace{1.8em}3.}{Weitere Versuche}{.\ .\ .\ .\ .\ .\ .\ .\ .\ .\ .\ .\ .\ .\ .\ .\ .\ .\ .\ .\ .\ .\ .\ .\ .\ .}{10}
\GCrow{\hspace{1.8em}4.}{\textsc{Gauss}’ Wiederholung des \textsc{Foucault}schen Versuches}{.\ .\ .\ .\ .\ .\ .\ .\ .\ .\ .\ .\ .}{12}
\GCrow{II.}{Das Prinzip des kleinsten Zwanges}{.\ .\ .\ .\ .\ .\ .\ .\ .\ .\ .\ .\ .\ .\ .\ .\ .\ .\ .\ .\ .}{16—32}
\GCrow{\hspace{1.8em}1.}{Historisches}{.\ .\ .\ .\ .\ .\ .\ .\ .\ .\ .\ .\ .\ .\ .\ .\ .\ .\ .\ .\ .\ .\ .\ .\ .\ .\ .\ .}{16}
\GCrow{\hspace{1.8em}2.}{Die \textsc{Gauss}sche Arbeit über das Prinzip des kleinsten Zwanges}{.\ .\ .\ .\ .\ .\ .\ .\ .}{21}
\GCrow{\hspace{1.8em}3.}{Das \textsc{Fourier}sche Ungleichungsprinzip}{.\ .\ .\ .\ .\ .\ .\ .\ .\ .\ .\ .\ .\ .\ .\ .\ .\ .}{26}
\GCrow{III.}{Die Anziehung der Ellipsoide}{.\ .\ .\ .\ .\ .\ .\ .\ .\ .\ .\ .\ .\ .\ .\ .\ .\ .\ .\ .\ .\ .\ .}{32—39}
\GCrow{\hspace{1.8em}1.}{\textsc{Gauss}’ Vorgänger}{.\ .\ .\ .\ .\ .\ .\ .\ .\ .\ .\ .\ .\ .\ .\ .\ .\ .\ .\ .\ .\ .\ .\ .\ .}{32}
\GCrow{\hspace{1.8em}2.}{Die \textsc{Gauss}sche Abhandlung}{.\ .\ .\ .\ .\ .\ .\ .\ .\ .\ .\ .\ .\ .\ .\ .\ .\ .\ .\ .\ .\ .}{36}
\GCrow{IV.}{Potentialtheorie}{.\ .\ .\ .\ .\ .\ .\ .\ .\ .\ .\ .\ .\ .\ .\ .\ .\ .\ .\ .\ .\ .\ .\ .\ .\ .\ .\ .}{40—52}
\GCrow{\hspace{1.8em}1.}{Die »Allgemeinen Lehrsätze«}{.\ .\ .\ .\ .\ .\ .\ .\ .\ .\ .\ .\ .\ .\ .\ .\ .\ .\ .\ .\ .\ .}{40}
\GCrow{\hspace{1.8em}2.}{Die Entwicklung bei \textsc{Gauss}}{.\ .\ .\ .\ .\ .\ .\ .\ .\ .\ .\ .\ .\ .\ .\ .\ .\ .\ .\ .\ .\ .}{47}
\GCrow{V.}{Arbeiten zur praktischen Mechanik}{.\ .\ .\ .\ .\ .\ .\ .\ .\ .\ .\ .\ .\ .\ .\ .\ .\ .\ .\ .\ .}{52—60}
\GCrow{\hspace{1.8em}1.}{Schwingungsdauer einer Magnetnadel}{.\ .\ .\ .\ .\ .\ .\ .\ .\ .\ .\ .\ .\ .\ .\ .\ .\ .}{52}
\GCrow{\hspace{1.8em}2.}{Torsionsbestimmungen}{.\ .\ .\ .\ .\ .\ .\ .\ .\ .\ .\ .\ .\ .\ .\ .\ .\ .\ .\ .\ .\ .\ .\ .}{56}
\GCrow{\hspace{1.8em}3.}{Bestimmung von Trägheitsmomenten}{.\ .\ .\ .\ .\ .\ .\ .\ .\ .\ .\ .\ .\ .\ .\ .\ .\ .\ .}{57}
\GCrow{\hspace{1.8em}4.}{Zur Theorie der Wage}{.\ .\ .\ .\ .\ .\ .\ .\ .\ .\ .\ .\ .\ .\ .\ .\ .\ .\ .\ .\ .\ .\ .\ .}{59}
\par\bigskip\centerline{\rule{.25\textwidth}{.3pt}}
\clearpage\endgroup
